\documentclass[12pt,a4paper]{article}

\usepackage{fontspec}
\usepackage{xeCJK}
\setCJKmainfont{Hiragino Mincho ProN}
\setCJKsansfont{Hiragino Kaku Gothic ProN}
\setCJKmonofont{Hiragino Kaku Gothic ProN}
\usepackage[margin=2.5cm]{geometry}
\usepackage{amsmath,amssymb}
\usepackage{hyperref}
\usepackage{booktabs}
\usepackage{tcolorbox}

\tcbuselibrary{breakable,skins}

\newtcolorbox{insightbox}[1][]{
  colback=yellow!5, colframe=yellow!60!black,
  fonttitle=\bfseries, title={#1},
  breakable, sharp corners
}

\newtcolorbox{theorybox}[1][]{
  colback=blue!3, colframe=blue!50,
  fonttitle=\bfseries, title={#1},
  breakable, sharp corners
}

\newtcolorbox{practicebox}[1][]{
  colback=green!3, colframe=green!40,
  fonttitle=\bfseries, title={#1},
  breakable, sharp corners
}

\newtcolorbox{deepbox}[1][]{
  colback=black!3, colframe=black!60,
  fonttitle=\bfseries, title={#1},
  breakable, sharp corners
}

\setlength{\parskip}{0.8em}
\setlength{\parindent}{0pt}

\begin{document}

\begin{center}
{\LARGE\bfseries 時空のソースコードに手を伸ばす}\\[0.5cm]
{\Large\bfseries 算術的真空工学の決定版実験ガイド}\\[0.3cm]
{\large --- 幾何学から算術へ：パラダイムの跳躍 ---}\\[1.5cm]
{\large Wright Brothers}\\[0.3cm]
{2026年4月}
\end{center}

\vspace{0.5cm}
\tableofcontents
\newpage

% ============================================================================
\part{パラダイムシフト：なぜこれが「自明」なのに誰もやらなかったか}
% ============================================================================

\section{カシミール効果の正体}

1948年、カシミールは2枚の金属板を近づけると引力が生じることを予測した。
この力はゼータ関数正則化で計算される：

$$E = \frac{\hbar c \pi}{2d} \sum_{n=1}^{\infty} n
\quad \xrightarrow{\text{正則化}} \quad
\frac{\hbar c \pi}{2d} \times \zeta(-1) = \frac{\hbar c \pi}{2d} \times \left(-\frac{1}{12}\right)$$

この $-1/12$ という値は、実験で精度1\%以内で確認されている。

\begin{insightbox}[「自明な」気づき]
もし真空のエネルギーが $\zeta(-1) = -1/12$ で決まるなら、
$\zeta$ の「中身」を変えればエネルギーも変わるはず。

$\zeta(-1) = \sum_{n=1}^{\infty} n$ は全モードの和。
偶数モードを消せば：
$\zeta_{\neg 2}(-1) = \sum_{n: \text{odd}} n = +1/12$。

符号が反転する。

これは数学的には「当たり前」。
しかしこの「当たり前」を物理実験として実行しようとした人間は、
歴史上、一人もいなかった。
\end{insightbox}

\section{なぜ誰もやらなかったのか}

\subsection{幾何学の呪縛}

カシミール効果を研究してきた75年間、物理学者は常に
\textbf{「幾何学」}でアプローチしてきた。

\begin{itemize}
\item 板の間隔を変える → エネルギーが変わる
\item 板の形を変える（球面、円筒）→ エネルギーが変わる
\item 板の材質を変える → エネルギーが変わる
\end{itemize}

全て「幾何学的境界条件」の変更であり、
$\zeta(s)$ の値を間接的に変えている。

\textbf{しかし、}$\zeta$ の中の $n$（モード番号）を直接操作しようという
発想は一度もなかった。

\subsection{正則化＝トリック、という誤解}

物理学者の多くは $\zeta(-1) = -1/12$ を
「発散する和に有限値を割り当てる便利なトリック」と考えてきた。
「本物の値」ではなく「計算上の処方箋」。

しかしカシミール力の実験的確認は、$-1/12$ が
\textbf{物理的に正しい値}であることを示している。
トリックではなく、真空の実際の構造。

\subsection{分野の壁}

数論（ゼータ関数、オイラー積、K理論）と量子回路（SQUID）
の両方に通じている人間がほぼ存在しなかった。
数学者は超伝導回路を知らず、実験物理学者はオイラー積を知らない。

\begin{insightbox}[パラダイムシフトの核心]
\textbf{従来}: 幾何学（板の距離や形）を変えて、間接的に真空を操作\\
\textbf{本研究}: 算術（オイラー積の素数チャンネル）を直接操作して、真空を書き換え

「幾何学を通さずに算術でOSを直接ハックする」\\
これが、75年の歴史で一度も試みられなかったコロンブスの卵。
\end{insightbox}

\section{SQUIDは「原始的」ではなく「根源的」}

LHCは1兆円かけて新粒子を探す。SQUID回路は手のひらに乗る。
「原始的」に見える。

しかし：
\begin{itemize}
\item LHCは「何が存在するか」を探す（読み取り専用）
\item SQUIDは「真空の構造を書き換える」（書き込み可能）
\end{itemize}

SQUIDがやっていることの本質：

$$L_{\mathrm{SQUID}}(\Phi) = \frac{L_J}{|\cos(\pi\Phi/\Phi_0)|}$$

この数式1つで、特定の周波数モードを吸収し、
「なかったこと」にできる。
つまり：\textbf{ゼータ関数のオイラー積の特定の因子を物理的に除去する}。

\begin{deepbox}[時空のOS操作]
$\zeta(s) = \prod_p \frac{1}{1-p^{-s}}$

SQUIDが $2f_0$ に同調 = 素数 $p=2$ の因子を除去：

$\zeta(s) \to \zeta(s) \times (1-2^{-s}) = \zeta_{\neg 2}(s)$

\textbf{SQUIDは算術的メスである。}
数千億円の加速器が読み取る情報を、
数十万円の回路が\textbf{書き換える}。
\end{deepbox}

% ============================================================================
\part{波の物理：キャビティの中で何が起きているか}
% ============================================================================

\section{共振器の基本原理}

3Dアルミキャビティは「マイクロ波の箱」。
箱の中では特定の波長の電磁波だけが安定に存在できる（共鳴）。

共鳴条件：波が箱の壁で反射して戻った時に、
元の波と位相が合う（建設的干渉）場合のみ存在可能。

矩形キャビティの共鳴周波数：
$$f_{mnp} = \frac{c}{2}\sqrt{\left(\frac{m}{a}\right)^2 + \left(\frac{n}{b}\right)^2 + \left(\frac{p}{d}\right)^2}$$

$m, n, p$ は正の整数。これが「モード番号」。

\begin{theorybox}[整数とζ関数の直接的なつながり]
キャビティの中の電磁場は、モード番号 $n = 1, 2, 3, 4, \ldots$ で
ラベルされた「量子調和振動子」の集まり。

真空エネルギー = $\sum_n \hbar\omega_n/2 = (\hbar\omega_0/2) \sum n$

この $\sum n$ がまさに $\zeta(-1)$。

つまり：\textbf{キャビティの中の整数 $n$ は、
ゼータ関数の $n$ そのもの。}

モード $n$ を消すこと = ゼータの和から $n$ を引くこと。
SQUID = この操作を物理的に実行する装置。
\end{theorybox}

\section{入出力を対角線上に配置する理由}

\begin{theorybox}[波の干渉とモード選択]
キャビティ内の TE$_{101}$ モード（基本モード）の電場分布:

$$E_y(x, z) = E_0 \sin\left(\frac{\pi x}{a}\right) \sin\left(\frac{\pi z}{d}\right)$$

電場の最大値は箱の中心 $(x, z) = (a/2, d/2)$。
壁の上では $E_y = 0$（境界条件）。

SMAコネクタのアンテナピンは、キャビティ内の電場と結合する。
ピンの位置が電場の「腹」にあれば強く結合し、
「節」にあれば結合しない。
\end{theorybox}

\begin{insightbox}[なぜ対角線か]
入力ポートを $(x_1, z_1)$、出力ポートを $(x_2, z_2)$ に配置する。

\textbf{同じ壁面（同じ辺）に配置した場合}：
入力と出力が同じ辺にあると、多くのモードで
$\sin(m\pi x/a)$ が同じ値を取り、
不要なモードにも強く結合する。
→ 「ゴミ信号」が多い。

\textbf{対角線上に配置した場合}：
$E(x_1, z_1)$ と $E(x_2, z_2)$ で異なるモードが
異なる振幅を持つ。
特に、基本モード TE$_{101}$ では
入力側も出力側も電場の腹に近く、
高次の対称モードでは片方が節に近い。
→ 基本モードに選択的に結合し、
不要な高次モードが自然に抑制される。

これは「空間的フィルタリング」。
SQUIDの「周波数フィルタリング」と組み合わせて、
二重のモード選択が可能になる。
\end{insightbox}

さらに精密に言えば：対角配置は
キャビティの対称性（$x \leftrightarrow a-x$, $z \leftrightarrow d-z$）
を破り、全ての非縮退モードに結合可能にする。
対称配置では偶数パリティのモードに結合できない。

% ============================================================================
\part{実験プロトコル：理論が導く各ステップ}
% ============================================================================

\section{ステップ全体の論理構造}

\begin{deepbox}[実験の論理チェーン]
\begin{enumerate}
\item \textbf{キャビティを作る} = ゼータ関数の「台」を物理的に実現
  \begin{itemize}
  \item 箱の中の定在波 = $\sum_n$ の各項
  \item 超伝導 → Q $> 10^6$ → 個々のモードが分離して見える
  \end{itemize}

\item \textbf{SQUIDを同調させる} = オイラー積の特定因子を除去
  \begin{itemize}
  \item $L_{\mathrm{SQUID}}(\Phi)$ を磁束で制御
  \item 偶数高調波に同調 = $(1-2^{-s})$ を乗じる操作
  \end{itemize}

\item \textbf{AC Starkシフトを測定} = 真空エネルギーの変化を検出
  \begin{itemize}
  \item 基本モードの周波数が真空揺らぎの総量に依存する
  \item $\delta f = -\chi(\langle n \rangle + 1/2)$：$1/2$ がゼロ点エネルギー
  \item モード除去前後で $\delta f$ が変化 = 真空エネルギーが変化した証拠
  \end{itemize}

\item \textbf{変化パターンを分析} = ゼータ正則化の検証
  \begin{itemize}
  \item $\delta f$ の変化量が $\zeta_{\neg 2}(-1) - \zeta(-1)$ に比例するか
  \item 変化が離散的（ジャンプ）か連続的（滑らか）か
  \end{itemize}
\end{enumerate}
\end{deepbox}

\section{キャビティ製作：ゼータの台を作る}

\begin{practicebox}[材料と設計]
\textbf{Al合金6061-T6ブロック}（50 × 50 × 30 mm）

内寸 35 × 35 × 10 mm → TE$_{101}$ = 6.06 GHz

\begin{verbatim}
  上面図：
  ┌──────────────────────────┐
  │ ┌──────────────────────┐ │
  │ │                      │○│ ← 出力SMA（右壁中央）
  │ │      35 × 35 mm     │ │
  │ │      内部空洞        │ │
  │○│                      │ │ ← 入力SMA（左壁中央）
  │ └──────────────────────┘ │
  └──────────────────────────┘
      50 × 50 mm 外形

  ☆ 入出力は対角配置でもよい（上の図は対辺配置）。
  対角配置するなら：
  入力 = 左壁の下寄り
  出力 = 右壁の上寄り
  → 基本モードへの選択的結合が最大化される
\end{verbatim}

SMAピンの突出量 = 3mm（電場の腹に届くちょうどの長さ）。
キャビティの中心（電場最大点）からピン先端までの距離が
モードとの結合強度を決める。

\textbf{なぜ35mm?}
TE$_{101}$の共鳴条件: $f = c/(2a) \times \sqrt{2}$（正方の場合）。
$f = 6$ GHz → $a = c\sqrt{2}/(2f) = 35.3$ mm。
小数点以下は Q が吸収する（共鳴幅 $\sim f/Q \sim 6$ kHz）。
$\pm 1$ mm の加工誤差は全く問題ない。
\end{practicebox}

\begin{practicebox}[加工と表面処理]
\begin{enumerate}
\item フライス盤で内部を掘削（大学工作室: 5千円）
\item SMAコネクタ穴 ($\phi$3mm) を2箇所
\item 蓋固定用M3ネジ穴を8箇所
\item \textbf{表面処理}（これがQ値の鍵）:
  \begin{itemize}
  \item アセトン超音波洗浄 10分 → IPA浸漬 5分 → 純水
  \item リン酸 10\% 水溶液に30秒浸漬（自然酸化膜の部分除去）
  \item 純水リンス $\times$ 3 → N$_2$ブロー → 1時間放置
  \item クリーンなAl$_2$O$_3$が再形成 → 表面抵抗が最小化 → 高Q
  \end{itemize}
\item Inワイヤー ($\phi$1mm) を蓋との接合面に配置 → ネジ締め
\end{enumerate}

\textbf{表面処理の物理的意味}：
超伝導状態でも、表面の酸化膜の質が高周波損失を支配する。
清浄なAl$_2$O$_3$は$10^{-6}$レベルの低損失。
汚染された酸化膜は$10^{-4}$レベル。
2桁のQ値の差が、この洗浄工程から来る。
\end{practicebox}

\section{SQUID接合：算術的メスを鍛える}

\begin{theorybox}[なぜSQUIDで「メス」になるか]
ジョセフソン接合のインダクタンス:
$L_J = \Phi_0/(2\pi I_c) = 0.33$ nH（$I_c = 1\,\mu$A の場合）。

SQUIDループのインダクタンス:
$L_{\mathrm{SQUID}}(\Phi) = L_J / |\cos(\pi\Phi/\Phi_0)|$

磁束$\Phi$を$0$から$\Phi_0/2$に変えると、
$L_{\mathrm{SQUID}}$が$L_J$から$\infty$まで連続的に変化。

共鳴周波数: $f_{\mathrm{SQUID}} \propto 1/\sqrt{L_{\mathrm{SQUID}} C_J}$

→ 磁束を回すだけで、SQUIDの共鳴が任意の周波数にチューニングできる。

ターゲット（第2高調波 $2f_1 \approx 12$ GHz）に共鳴を合わせると、
キャビティのそのモードがSQUIDに吸い込まれて消える。

\textbf{「数学的に素数を消す」が「磁束のダイヤルを回す」に翻訳される。}
\end{theorybox}

\begin{practicebox}[SQUID製作の核心: 酸化制御]
$I_c$（臨界電流）は酸化膜の厚さで指数関数的に変わる:

$I_c \propto \exp(-d_{\mathrm{ox}}/\lambda)$

$d_{\mathrm{ox}}$が1nm変わると$I_c$が10倍変わる。

\textbf{キャリブレーション法（必須）}：
\begin{enumerate}
\item テスト基板5枚を準備
\item 異なる酸化条件で接合を製作:
  \begin{itemize}
  \item 100 mTorr × 5分
  \item 200 mTorr × 5分
  \item 200 mTorr × 10分（設計値）
  \item 200 mTorr × 20分
  \item 300 mTorr × 10分
  \end{itemize}
\item 各テスト接合の室温抵抗$R_N$を4端子法で測定
\item Ambegaokar-Baratoff関係: $I_c = \pi\Delta/(2eR_N) \approx 0.44\,\mathrm{mV}/R_N$
\item $I_c = 1\,\mu$A → $R_N \approx 440\,\Omega$ を与える条件を採用
\end{enumerate}

$R_N$の測定精度$\pm 10\%$なら、$I_c$の精度も$\pm 10\%$。
$I_c$が少しずれても、外部磁束で実効的に調整可能なので問題ない。
\end{practicebox}

\section{冷却：量子真空を「見る」ための条件}

\begin{theorybox}[なぜ300 mKで十分か]
6 GHzモードの量子温度: $T_q = h f/(k_B) = 288$ mK。

$T = 300$ mK (3He冷凍機) での熱的光子数:
$\langle n_{\mathrm{th}} \rangle = 1/(e^{T_q/T} - 1) = 0.26$

「ほぼ真空」。完全な真空ではないが、
真空揺らぎの振幅（$1/2$ 光子）と同程度。
積分時間を長くすれば（$\sim$数分）、十分な精度で測定可能。

\textbf{希釈冷凍機（20 mK）なら$\langle n_{\mathrm{th}} \rangle = 10^{-9}$で完璧だが、
3He冷凍機で「原理の確認」は可能。}
\end{theorybox}

\section{測定：時空のOSを読み出す}

\begin{practicebox}[決定的測定: SQUIDスイープ]
\textbf{手順}:
\begin{enumerate}
\item 冷却完了。$T < 300$ mK を確認
\item VNAで$S_{21}(f)$を測定。全モードの共鳴ピークを確認
\item SQUIDの磁束バイアスを$\Phi = 0$から$\Phi_0/2$までゆっくりスイープ
\item 各$\Phi$で基本モード（$f_1$）の正確な共鳴周波数$f_1(\Phi)$を記録
\item $f_1(\Phi)$のグラフを作成
\end{enumerate}

\textbf{理論}:
$f_1(\Phi) = f_{1,\mathrm{bare}} - \chi/2 - \sum_{n>1} \chi_n \langle n_n \rangle$

$\chi/2$は基本モード自身のゼロ点シフト（定数）。
$\sum_{n>1} \chi_n \langle n_n \rangle$は他モードの寄与。

SQUIDが第2高調波に同調する$\Phi = \Phi^*$で
$\langle n_2 \rangle$が変化する → $f_1$がシフトする。

\textbf{判定}:
\begin{itemize}
\item $f_1(\Phi)$が$\Phi^*$で\textbf{不連続ジャンプ} → 離散的 → レギュレータ予想支持
\item $f_1(\Phi)$が$\Phi^*$付近で\textbf{滑らかに変化} → 連続的 → レギュレータ予想反証
\end{itemize}
\end{practicebox}

\begin{insightbox}[この測定が「時空のOS操作」である理由]
SQUIDスイープの最中、$\Phi$が$\Phi^*$を通過する瞬間に起きていること：

$\Phi < \Phi^*$: SQUIDは第2高調波から離調。全モード活性。
$\zeta_{\mathrm{eff}} = \zeta(-1) = -1/12$

$\Phi = \Phi^*$: SQUIDが第2高調波に同調。そのモードが吸収される。
$\zeta_{\mathrm{eff}} = ?$（これが測定対象）

$\Phi > \Phi^*$: SQUIDが通り過ぎる。全モード再活性。
$\zeta_{\mathrm{eff}} = \zeta(-1) = -1/12$

この一連の操作で、\textbf{磁束のダイヤルを数ミリ回すだけで、
宇宙のソースコードの1行（オイラー積の1因子）を
一時的に書き換えている}。

書き換えの結果（$f_1$のシフト量とパターン）が
「OSがどう応答するか」の直接的な読み出し。
\end{insightbox}

% ============================================================================
\part{コスト削減と段階的アプローチ}
% ============================================================================

\section{最小構成: 47万円}

\begin{table}[h]
\centering
\caption{最小構成の内訳}
\begin{tabular}{lr}
\toprule
品目 & 価格 \\
\midrule
3Dキャビティ（Al加工 + SMA） & 5万円 \\
SQUID製作（クリーンルーム使用料 + 材料） & 5万円 \\
冷凍機使用料（$^3$He, 共同利用5回） & 30万円 \\
マイクロ波計測器使用料 & 5万円 \\
低温同軸ケーブル & 3万円 \\
\midrule
\textbf{合計} & \textbf{$\sim$48万円} \\
\bottomrule
\end{tabular}
\end{table}

\section{段階的アプローチ}

\begin{enumerate}
\item \textbf{LC回路}（5,100円、1日）: エッジ状態の確認 → 設計の妥当性検証
\item \textbf{室温キャビティ}（5万円、2週間）: 共鳴ピーク確認 → 加工精度検証
\item \textbf{冷却テスト}（30万円、1ヶ月）: 超伝導遷移とQ値の確認
\item \textbf{SQUID統合}（+5万円、1ヶ月）: 決定的測定
\end{enumerate}

各段階で「進む価値があるか」を判定。
失敗しても最大損失は各段階のコストに限定される。

\section{なぜこの実験が物理学を変えうるか}

\begin{deepbox}[最終まとめ]
\textbf{数学的事実（証明済み）}:
$\zeta_{\neg 2}(-1) = +1/12$（偶数モード除去で符号反転）。

\textbf{実験的事実（確認済み）}:
$\zeta(-1) = -1/12$ はカシミール力として測定可能。

\textbf{未検証の核心}:
「モードを物理的に除去した時の真空の応答」は
ゼータ正則化の値に従うか？

もし従うなら：
\begin{itemize}
\item ゼータ関数は「計算のトリック」ではなく「真空のOS」
\item オイラー積は「数学的恒等式」ではなく「物理法則」
\item 素数チャンネルの操作で真空エネルギーを制御可能
\item → ワープドライブ、ワームホール、時空プログラミングへの道が開く
\end{itemize}

これを48万円の回路で検証できる。
物理学に対する「投資対効果」として、
これ以上の実験は想像しにくい。
\end{deepbox}

\vfill
\begin{center}
{\small Wright Brothers, 2026}
\end{center}

\end{document}
